\section{Salience (Prevalence)}

Using the Trend Track, yearly mention rates will be estimated as relative frequencies (candidate and validated hits per scanned corpus size), enabling a diachronic test of whether these terms are becoming more prominent in general web discourse. \textit{[status: incomplete]}

\section{Intensity (Vertical Creep)}

Collocates within a $\pm 5$-word window of the target term will be extracted per year and scored using the NRC Valence--Arousal--Dominance lexicon (M. Saif, 2025). Annual weighted arousal indices (and derived ``severity'' composites, where appropriate) will test whether the surrounding language becomes more or less emotionally intense over time. \textit{[status: incomplete]}

\section{Breadth (Horizontal Creep)}

Breadth will be estimated via contextual dispersion: a fixed number of sentences containing the target term will be sampled per year, embedded with transformer sentence embeddings, and summarised by average pairwise cosine distance. Increasing dispersion indicates a widening range of contexts (semantic broadening). \textit{[status: incomplete]}

\section{Sentiment (Connotation)}

Using the same collocate windows, annual weighted valence indices (NRC VAD) will test for shifts toward more positive vs more negative connotations (amelioration vs degeneration), reflecting potential destigmatisation or increased pejoration. \textit{[status: incomplete]}

\section{Thematic Evolution}

The Corpus Track will be analysed with BERTopic (transformer embeddings, UMAP + HDBSCAN clustering, c-TF-IDF representations) to trace topics-over-time in annual bins. \textit{[status: incomplete]}

\section{Analytic Strategy}

Linear regression analyses were performed to test the statistical significance of the predicted trends in the hypotheses....
